\chapter{Nonlinear SSM Validation Ladder}
\label{ch:bf-nonlinear-ssm-validation}

The nonlinear SSM ladder tests whether approximate nonlinear filtering can be
used inside HMC before DSGE perturbation solvers and economic boundaries enter
the picture.  It should use small synthetic models with known parameters and
independent references, not large application models.

\section{Controlled Models}
\label{sec:bf-nonlinear-controlled-models}

The first nonlinear validation suite is deliberately small.  It is designed to
test the filtering law, structural deterministic completion, and sigma-point
moment approximation before DSGE perturbation solvers, client adapters, or HMC
geometry enter the picture.  The suite contains three first-rung models.

\subsection{Model A: Affine Gaussian Structural Oracle}
\label{subsec:bf-nonlinear-model-a}

The affine oracle is a BayesFilter-local synthetic model:
\[
  x_t=(m_t,\ell_t)^\top,\qquad
  m_t=0.35m_{t-1}-0.10\ell_{t-1}+0.25\varepsilon_t,\qquad
  \ell_t=m_{t-1},
\]
\[
  y_t=m_t+\eta_t,\qquad
  \varepsilon_t\sim N(0,1),\qquad
  \eta_t\sim N(0,0.15^2).
\]
The initial law is Gaussian and the deterministic coordinate \(\ell_t\)
stores the previous \(m\).  The oracle is the exact linear Gaussian Kalman
likelihood after converting the structural law to its full-state LGSSM.  The
test must verify that SVD cubature, SVD-UKF, and SVD-CUT4 collapse to the same
prediction-error likelihood and that deterministic-completion residuals are
zero up to numerical tolerance.

\subsection{Model B: Nonlinear Accumulation}
\label{subsec:bf-nonlinear-model-b}

The nonlinear accumulation model is a smooth BayesFilter-local structural
fixture:
\[
  m_t=\rho m_{t-1}+\sigma\varepsilon_t,\qquad
  k_t=\alpha k_{t-1}+\beta\tanh(m_t),
\]
\[
  y_t=m_t+k_t+\eta_t,\qquad
  \varepsilon_t\sim N(0,1),\qquad
  \eta_t\sim N(0,\sigma_\eta^2).
\]
The first V1 defaults are
\[
  \rho=0.70,\qquad \sigma=0.25,\qquad \alpha=0.55,\qquad
  \beta=0.80,\qquad \sigma_\eta=0.30.
\]
This model tests nonlinear deterministic completion without a time-indexed
transition hook.  The reference oracle is a dense one-step Gaussian
moment-projection quadrature: it integrates the augmented predictive Gaussian,
computes the transformed prediction and observation moments, and then applies
the same Gaussian update used by the sigma-point filters.  It is not an exact
nonlinear likelihood.

\subsection{Model C: Nonlinear Growth Benchmark}
\label{subsec:bf-nonlinear-model-c}

The scalar nonlinear growth benchmark is the standard model associated with
Kitagawa's Monte Carlo filtering examples \citep{kitagawa1996montecarlo} and
widely reused in sequential Monte Carlo discussions.  Its time-inhomogeneous
law is
\[
  x_t=\frac{x_{t-1}}{2}
      +\frac{25x_{t-1}}{1+x_{t-1}^2}
      +8\cos(1.2(t-1))
      +\sigma_x\varepsilon_t,
\]
\[
  y_t=\frac{x_t^2}{20}+\sigma_y\eta_t,
  \qquad \varepsilon_t,\eta_t\sim N(0,1).
\]
Because the current BayesFilter structural TensorFlow transition does not pass
an explicit time index, the V1 testing fixture uses an autonomous deterministic
phase coordinate:
\[
  \tau_t=\tau_{t-1}+1,\qquad \tau_1=1,
\]
\[
  x_t=\frac{x_{t-1}}{2}
      +\frac{25x_{t-1}}{1+x_{t-1}^2}
      +8\cos(1.2\tau_{t-1})
      +\sigma_x\varepsilon_t.
\]
The observation equation remains \(y_t=x_t^2/20+\sigma_y\eta_t\).  The first
defaults are \(x_1\sim N(0,0.2)\), \(\sigma_x=1\), and \(\sigma_y=1\).
As in Model B, the first oracle is dense one-step Gaussian moment projection,
not the exact nonlinear likelihood.

\subsection{Deferred Models}
\label{subsec:bf-nonlinear-deferred-models}

\begin{description}
  \item[Bearings-only tracking.]
    \citet{gordon1993novel} and \citet{arulampalam2002tutorial} are useful
    sources, but the test should wait until angle residual and wrapping policy
    are fixed.
  \item[Radar range-bearing tracking.]
    This should follow the bearings-only residual policy and should add
    mixed-scale observation diagnostics.
  \item[Stochastic volatility.]
    This should wait until the sigma-point interface supports non-additive
    observation noise or an explicit approximation contract, since current
    filters assume additive observation covariance after \(h(x_t)\).
\end{description}

Deferred means that the model is useful later, not that it is rejected.  The
first V1 suite stops at Models A--C so that value, score, branch, and benchmark
tests share a stable set of fixtures.

\section{Required Outputs}
\label{sec:bf-nonlinear-required-outputs}

A nonlinear SSM HMC validation run should write:
\begin{itemize}
  \item JSON diagnostics with chain counts, warmup/draw counts, step sizes,
    acceptance rate, divergence counts, R-hat, ESS, MCSE/SD, and runtime;
  \item NPZ chain artifacts for independent audit;
  \item true-parameter coverage when synthetic truth exists;
  \item filter failure counts, conditioning-guard counts, and finite-gradient
    failures.
\end{itemize}
Medium runs are sampler-usability evidence.  Strict runs require at least four
chains, finite log targets, zero or justified divergences, R-hat below 1.01,
bulk and tail ESS above the planned thresholds, and saved artifacts.

\section{Stop Rules}
\label{sec:bf-nonlinear-stop-rules}

If the LGSSM rung fails, the nonlinear rung should not start.  If the nonlinear
SSM rung fails while LGSSM passes, the failure should be attributed first to
the approximation, nonlinear target geometry, transition hook, or derivative
path.  DSGE/NK debugging should wait until this simpler explanation is ruled
out.
